\documentclass{article}
\usepackage[T2A]{fontenc}
\usepackage[utf8]{inputenc}
\usepackage{amsthm}
\usepackage{amsmath}
\usepackage{amssymb}
\usepackage{amsfonts}
\usepackage{mathrsfs}
\usepackage[12pt]{extsizes}
\usepackage{fancyvrb}
\usepackage{indentfirst}
\usepackage[
  left=2cm, right=2cm, top=2cm, bottom=2cm, headsep=0.2cm, footskip=0.6cm, bindingoffset=0cm
]{geometry}
\usepackage[english,russian]{babel}


\begin{document}
\section*{Вариант 10}

Решение уравнений с учетом краевых условий представим в виде

\begin{equation}
    W = \sum\limits_{k=1}^{\infty} \left[\left(R_k(\tau)+R^0_k\right) \cos{\frac{(2k-1)\pi\xi}{2} + Q_k(\tau) \sin{k\pi\xi}} \right]
\end{equation}
Верхний индекс 0 в (1) означает решение, соответствующее постоянному уровню давления $p_0$, независящему от $\tau$.

Принимая во внимание линейность предыдущего уравнения и подставляя в него (1), найденное выражение для давления, а, также раскладывая оставшиеся члены, входящие в его правую часть в ряды по тригонометрическим функциям, из полученного уравнения запишем выражение для составляющей, независящей от времени $R^0_k$
\begin{equation*}
    R^0_k = \left(2\ell/\left(\left(2k-1\right)\pi\right)\right)^4 \left(4\left(-1\right)^{k+1}/\left(\left(2k-1\right)\pi\right)\right) p_0\left(Dw_m\right)^{-1},
\end{equation*}
и уравнения для определения $R_k(\tau)$ и $Q_k(\tau)$
\begin{equation}
    a_{1ck}w_m R_k +  a_{2ck}w_m d R_k / d\tau = 2(-1)^{k+1}/((2k-1)\pi D) p_m^+ f_p(\tau),
\end{equation}
\begin{equation}
    a_{1sk}w_m Q_k + a_{2sk} w_m Q_k/d\tau = (-1)^{k+1}/(k\pi D)p_m^+ f_p(\tau).
\end{equation}
\end{document}
